\assignq{4}{3}{Neutral sodium: radial eigenstates, E1 rates and lifetimes}{1}
This is the computational question. Here we solve
a stated effective radial model numerically and
show all E1 rate calculations; the corresponding
reproducible Python script accompanies the booklet
at \texttt{scripts/assignment\_numerical\_solutions.py}.

\subsection*{Step 1: choose the effective central potential}
Use the parametric screened potential previously
considered in the sodium practical:
\[
\boxed{V_b(r)=-\frac Zr+
\frac{Z-\xi}{r}(1-e^{-r/b}),\qquad
Z=11,\ \xi=1,\ b=0.4072},
\]
with \(r,b\) in atomic units. Equivalently
\(V=-1/r-10e^{-r/b}/r\).
It approaches \(-11/r\) close to the nucleus
and \(-1/r\) in the valence region.
It is a \emph{single-electron effective model},
not a full electron--electron Hartree--Fock
calculation.

For \(u_{nl}(r)=rR_{nl}(r)\), solve
\[
\boxed{-\frac12u''(r)+
\left[\frac{l(l+1)}{2r^2}+V_b(r)\right]u(r)
=E_{nl}u(r)},
\quad
u(0)=u(r_{\max})=0,\quad
\int_0^{r_{\max}}|u|^2\diff r=1.
\]
Use a finite-difference grid of spacing
\(h=0.012\,a_0\) and \(r_{\max}=120\,a_0\).
The second derivative uses
\(u''(r_i)\simeq(u_{i-1}-2u_i+u_{i+1})/h^2\).
The resulting Hamiltonian is tridiagonal:
\[
H_{ii}=\frac1{h^2}+\frac{l(l+1)}{2r_i^2}+V(r_i),
\quad H_{i,i\pm1}=-\frac1{2h^2}.
\]
Select the correct \(n,l\) eigenstates by
radial node index \(n-l-1\): \(3s\) has two
nodes, \(4s\) three, \(3p\) one, \(3d\) zero.

\subsection*{Step 2: numerical energy levels}
The calculated energies (relative to
ionisation, in Hartree) are
\[
\begin{array}{c|r|r}
\text{State}&l&E_{nl}\ (\mathrm{Ha})\\\hline
3s&0&-0.1831367\\
3p&1&-0.1102218\\
4s&0&-0.0705131\\
3d&2&-0.0556813
\end{array}
\]
The experimentally compiled sodium excitation
energies for comparison are approximately
\(16956\ \mathrm{cm}^{-1}\) for \(3p_{1/2}\),
\(25740\ \mathrm{cm}^{-1}\) for \(4s\), and
\(29173\ \mathrm{cm}^{-1}\) for \(3d\).
They are not exactly reproduced because
a single spin-independent central potential
cannot capture all fine-structure and
correlation effects.

\subsection*{Step 3: identify every downward E1 channel}
For a single electron, the E1 dipole operator
has odd parity and \(\Delta l=\pm1\).
With the ordering above, the primary allowed
decays are
\[
\boxed{3p\to3s,\qquad
4s\to3p,\qquad
3d\to3p.}
\]
The direct \(4s\to3s\) and \(3d\to3s\)
channels are E1-forbidden (\(\Delta l=0\)
and \(\Delta l=-2\), respectively).
The model level sketch is:
\[
\begin{tikzpicture}[x=1cm,y=.9cm,>=Latex]
\draw[->] (0,0)--(0,5.1) node[above]{Energy};
\draw[accent,very thick](1,0.7)--(3.4,0.7)
node[right,ink]{\(3s\;{}^2S\)};
\draw[accent,very thick](1,2.4)--(3.4,2.4)
node[right,ink]{\(3p\;{}^2P^o\)};
\draw[accent,very thick](1,3.4)--(3.4,3.4)
node[right,ink]{\(4s\;{}^2S\)};
\draw[accent,very thick](1,4.2)--(3.4,4.2)
node[right,ink]{\(3d\;{}^2D\)};
\draw[->,thick] (1.5,2.25)--(1.5,.83);
\draw[->,thick] (2.1,3.25)--(2.1,2.53);
\draw[->,thick] (2.8,4.05)--(2.8,2.53);
\end{tikzpicture}
\]

\subsection*{Step 4: evaluate the radial dipole integrals}
With \(u=rR\), the radial matrix element is
\[
D_{if}=\int_0^\infty u_f(r)\,r\,u_i(r)\diff r.
\]
The numerically computed magnitudes and
transition energies are
\[
\begin{array}{c|r|r}
i\to f&|D_{if}|\ (a_0)&
\omega_{if}=E_i-E_f\ (\mathrm{Ha})\\\hline
3p\to3s&4.48268&0.0729149\\
4s\to3p&4.36430&0.0397087\\
3d\to3p&5.97665&0.0545405
\end{array}
\]
Angular integration over \(Y_{lm}\) supplies
the E1 weights. For a nonrelativistic
spinless orbital with its magnetic substates
appropriately averaged, the total downward
rate is
\[
\boxed{A_{if}=\frac{4\alpha_{\rm fs}^3}{3t_{\rm au}}\,
\omega_{if}^3\,\frac{l_>}{2l_i+1}\,
|D_{if}|^2,}
\]
where \(l_>=\max(l_i,l_f)\),
\(t_{\rm au}=2.4188843266\times10^{-17}\) s.
The factors \(l_>/(2l_i+1)\) are
\(1/3\) for \(p\to s\), \(1\) for
\(s\to p\) and \(2/5\) for \(d\to p\).

\subsection*{Step 5: sum rates and calculate lifetimes}
The excited-state lifetime is
\(\tau_i=[\sum_{f<i}A_{if}]^{-1}\).
With the specified model, each upper state has
one primary orbital E1 channel:
\[
\begin{array}{c|r|r}
i&\text{Calculated }A_i\ (\mathrm{s}^{-1})
&\boxed{\tau_i\ (\mathrm{ns})}\\\hline
3p&5.562\times10^7&17.979\\
4s&2.545\times10^7&39.147\\
3d&4.965\times10^7&20.139
\end{array}
\]
These are independently computed predictions,
not values fitted to the reference lifetimes.

\subsection*{Step 6: compare with published atomic data}
Published nonrelativistic/relativistic
atomic-structure calculations for the Na-like
sequence at \(Z=11\) report approximately
\(16.48\) ns for \(3p\), \(38.25\) ns for
\(4s\), and \(19.67\) ns for \(3d\);
the \(J\)-resolved values vary slightly.
Our percentage differences from these
\emph{theoretical reference} results are:
\[
\begin{array}{c|r|r|r}
&\text{this model (ns)}&
\text{ref. theory (ns)}&
\text{difference}\\\hline
3p&17.979&16.48&+9.10\%\\
4s&39.147&38.25&+2.34\%\\
3d&20.139&19.67&+2.39\%
\end{array}
\]
A separately published direct measurement
of Na I \(3p\) reports \(17.3\pm1.0\) ns,
consistent with our predicted \(17.98\) ns
at that uncertainty. The \(4s\) and \(3d\)
benchmark values in the table are
\emph{calculations}, not claimed direct
measurements; a precise experimental
comparison should use a cited measured
dataset for each fine-structure component.

\subsection*{Step 7: limitations}
This simple model ignores core polarisation,
electron correlation beyond the fitted
potential, spin--orbit fine structure and
configuration mixing. Its chosen \(b\)
parameter controls the radial overlaps;
different fitted potentials change \(\tau\).
Higher multipoles, collisional quenching,
cascades and branching into more levels
may matter outside the minimal model.
Numerical convergence should be checked by
decreasing \(h\) and increasing \(r_{\max}\).

\methodnote{Reproduce every number using
\texttt{python scripts/assignment\_numerical\_solutions.py}.
References: Froese Fischer \emph{et al.},
Atomic Data and Nuclear Data Tables \textbf{92}
(2006), Na-like \(Z=11\) lifetimes;
NIST Na I energy-level table; 1970 JOSA
Na I \(3p\) lifetime measurement.}
